%%================================================
%% Chapter 2: Theory and Related Work
%%================================================

\chapter{ทฤษฎีและงานวิจัยที่เกี่ยวข้อง}
\label{chapter2}

บทนี้รวบรวมแนวคิด ทฤษฎี และงานวิจัยที่เกี่ยวข้องกับระบบดูแลผู้สูงอายุและผู้ใช้รถเข็นในสภาพแวดล้อมอัจฉริยะ โดยเน้นเฉพาะพื้นฐานที่จำเป็นต่อการออกแบบระบบ ได้แก่ แนวคิด Ambient Assisted Living, เทคโนโลยี IoT และเซ็นเซอร์, การระบุตำแหน่งภายในอาคารด้วย BLE RSSI, การประมวลผลข้อมูล IMU, Machine Learning สำหรับจำแนกพฤติกรรม, Retrieval-Augmented Generation และแนวคิด Neuro-symbolic AI แบบผสมระหว่างกฎกับโมเดลภาษา

\section{แนวคิดระบบดูแลผู้สูงอายุและผู้ใช้รถเข็น}
\label{sec:ch2_elderly_wheelchair_care}

การดูแลผู้สูงอายุและผู้ใช้รถเข็นเป็นงานที่ต้องติดตามหลายมิติพร้อมกัน ทั้งตำแหน่งของผู้พักอาศัย การเคลื่อนไหว ความผิดปกติของสัญญาณชีพ การตอบสนองต่อเหตุฉุกเฉิน และภาระงานของบุคลากรดูแล ในบริบทของสังคมสูงวัย จำนวนผู้สูงอายุที่ต้องพึ่งพาผู้ดูแลมีแนวโน้มเพิ่มขึ้นต่อเนื่อง ขณะที่บุคลากรในสถานดูแลมีข้อจำกัดด้านเวลาและจำนวนคน การออกแบบระบบช่วยดูแลจึงควรช่วยลดงานตรวจสอบซ้ำ ๆ ของผู้ดูแล และทำให้ข้อมูลสำคัญถูกรวมอยู่ในจุดที่สามารถมองเห็นและตัดสินใจได้เร็วขึ้น \cite{WHO2021DecadeAgeing,Schulz2020CaregiverBurden}

แนวคิด Ambient Assisted Living (AAL) เป็นแนวคิดที่ใช้เทคโนโลยีสารสนเทศและสภาพแวดล้อมอัจฉริยะเพื่อสนับสนุนการใช้ชีวิตประจำวันของผู้สูงอายุหรือผู้ที่มีข้อจำกัดด้านร่างกาย ระบบในกลุ่มนี้มักประกอบด้วยเซ็นเซอร์ อุปกรณ์สื่อสาร ฐานข้อมูล และกลไกแจ้งเตือน เพื่อช่วยให้ผู้ดูแลรับรู้เหตุการณ์สำคัญได้ทันเวลาโดยไม่ต้องตรวจด้วยตนเองตลอดเวลา \cite{Rashidi2013AALSurvey} ในงานด้านผู้ใช้รถเข็น ข้อมูลตำแหน่งและการเคลื่อนไหวมีความสำคัญมากเป็นพิเศษ เพราะสามารถสะท้อนรูปแบบกิจกรรม ความเสี่ยงในการหยุดนิ่งนานผิดปกติ หรือการเคลื่อนที่ออกจากบริเวณที่กำหนด

Smart environment เป็นการขยายแนวคิด AAL ไปสู่ระดับพื้นที่ โดยทำให้อุปกรณ์และบริการภายในอาคารเชื่อมโยงกัน เช่น เซ็นเซอร์ประจำตัวผู้ใช้ อุปกรณ์สวมใส่ จุดอ้างอิงสัญญาณภายในอาคาร ระบบแจ้งเตือน และ dashboard สำหรับผู้ดูแล จุดสำคัญไม่ใช่เพียงการเก็บข้อมูลให้ได้มากที่สุด แต่เป็นการจัดข้อมูลให้มีบริบท เช่น ใครอยู่ที่ไหน กำลังเคลื่อนไหวหรือหยุดนิ่ง มีสัญญาณชีพผิดปกติหรือไม่ และควรส่งต่อเหตุการณ์ใดให้ผู้ดูแลตรวจสอบก่อน

จากแนวคิดดังกล่าว ระบบดูแลผู้สูงอายุและผู้ใช้รถเข็นจึงควรมีคุณสมบัติหลัก 4 ประการ ได้แก่ การรับข้อมูลจากเซ็นเซอร์ที่สัมพันธ์กับสถานะจริงของผู้ใช้ การส่งข้อมูลเข้าสู่ระบบกลางอย่างต่อเนื่อง การจัดเก็บและแสดงผลข้อมูลในรูปแบบที่ผู้ดูแลอ่านได้เร็ว และการมีกลไกช่วยจัดลำดับความสำคัญของเหตุการณ์ก่อนส่งให้มนุษย์ตรวจสอบ ทั้งหมดนี้เป็นพื้นฐานของการออกแบบระบบต้นแบบในงานนี้ แต่ในบทนี้จะอธิบายในระดับทฤษฎีเพื่อไม่ให้ซ้ำกับรายละเอียดการพัฒนาในบทถัดไป

\section{เทคโนโลยี IoT และ Sensor สำหรับระบบดูแล}
\label{sec:ch2_iot_sensor}

Internet of Things (IoT) หมายถึงการเชื่อมต่ออุปกรณ์กายภาพเข้ากับระบบสารสนเทศเพื่อเก็บ ส่ง และประมวลผลข้อมูลจากสภาพแวดล้อมจริง อุปกรณ์ IoT ในระบบดูแลผู้สูงอายุอาจเป็นเซ็นเซอร์ติดรถเข็น อุปกรณ์สวมใส่ โทรศัพท์มือถือ หรือ gateway ภายในพื้นที่ทดลอง จุดร่วมของอุปกรณ์เหล่านี้คือมีข้อจำกัดด้านพลังงาน ขนาด หน่วยความจำ และคุณภาพสัญญาณ จึงต้องเลือก protocol และโครงสร้างข้อมูลให้เหมาะกับงาน \cite{Atzori2010IoTSurvey}

Bluetooth Low Energy (BLE) เป็นเทคโนโลยีสื่อสารระยะใกล้ที่เหมาะกับอุปกรณ์ใช้พลังงานต่ำ BLE รองรับทั้งการโฆษณาแพ็กเก็ตแบบ advertising และการเชื่อมต่อผ่าน Generic Attribute Profile (GATT) ทำให้อุปกรณ์สามารถส่งข้อมูลสั้น ๆ เช่น RSSI, battery, IMU หรือค่าจาก sensor ไปยัง gateway ได้ โดยค่า Received Signal Strength Indicator (RSSI) สามารถนำไปใช้เป็นข้อมูลประกอบการประเมินระยะหรือห้องที่อุปกรณ์อยู่ได้ แม้ค่า RSSI จะมีความผันผวนจากสิ่งกีดขวาง การสะท้อนของสัญญาณ และตำแหน่งของอุปกรณ์ \cite{BluetoothSIGLEPrimer}

\begin{figure}[htbp]
    \centering
    \resizebox{0.92\textwidth}{!}{%
    \begin{tikzpicture}[
        font=\small,
        box/.style={draw, thick, align=center, text width=2.8cm, minimum height=1.0cm},
        arrow/.style={-{Latex[length=2mm]}, thick},
        dasharrow/.style={-{Latex[length=2mm]}, thick, dashed}
    ]
        \node[box, fill=blue!8] (sensor) {BLE Sensor\\(Peripheral)};
        \node[box, fill=gray!10, right=1.0cm of sensor] (adv) {Advertising\\Packet (RSSI)};
        \node[box, fill=green!8, right=1.0cm of adv] (scan) {Scan / Observe\\(Central)};
        \node[box, fill=orange!10, right=1.0cm of scan] (gatt) {GATT Connect\\อ่าน IMU, HR};
        \node[box, fill=purple!8, right=1.0cm of gatt] (gw) {Mobile\\Gateway};
        \draw[arrow] (sensor) -- (adv);
        \draw[arrow] (adv) -- (scan);
        \draw[dasharrow] (scan) -- node[above, font=\scriptsize]{เชื่อมต่อ} (gatt);
        \draw[arrow] (gatt) -- (gw);
    \end{tikzpicture}}
    \caption{กระบวนการสื่อสาร BLE: advertising สำหรับ RSSI และ GATT สำหรับข้อมูล sensor}
    \label{fig:ch2_ble_advertising_flow}
\end{figure}


MQTT เป็น protocol แบบ publish/subscribe ที่นิยมในระบบ IoT เพราะ payload มีขนาดเล็กและแยกผู้ส่งกับผู้รับออกจากกัน อุปกรณ์หรือ gateway สามารถ publish telemetry ไปยัง topic ที่กำหนด ขณะที่ server หรือ dashboard subscribe เฉพาะ topic ที่ต้องใช้ หลักการนี้เหมาะกับระบบที่มีข้อมูลหลายชนิด เช่น ตำแหน่ง สัญญาณชีพ และสถานะอุปกรณ์ เพราะทำให้เพิ่มแหล่งข้อมูลใหม่ได้โดยไม่ต้องผูกทุก service เข้าหากันโดยตรง \cite{MQTTSpec,Mishra2018_MQTT_survey}

\begin{figure}[htbp]
    \centering
    \resizebox{0.92\textwidth}{!}{%
    \begin{tikzpicture}[
        font=\small,
        box/.style={draw, thick, align=center, text width=2.6cm, minimum height=1.0cm},
        broker/.style={draw, thick, align=center, text width=2.6cm, minimum height=1.3cm, fill=yellow!12},
        arrow/.style={-{Latex[length=2mm]}, thick}
    ]
        \node[box, fill=blue!8] (wc) {Wheelchair\\Sensor};
        \node[box, fill=blue!8, below=0.5cm of wc] (wear) {Wearable\\Sensor};
        \node[box, fill=green!8, right=0.8cm of wc, yshift=-0.5cm] (gw) {Mobile\\Gateway};
        \node[broker, right=0.9cm of gw] (brk) {MQTT\\Broker};
        \node[box, fill=orange!10, right=0.9cm of brk, yshift=0.45cm] (api) {API\\Server};
        \node[box, fill=purple!8, right=0.9cm of brk, yshift=-0.45cm] (dash) {Dashboard};
        \draw[arrow] (wc) -- (gw);
        \draw[arrow] (wear) -- (gw);
        \draw[arrow] (gw) -- node[above, font=\scriptsize]{publish} (brk);
        \draw[arrow] (brk) -- node[above, font=\scriptsize]{subscribe} (api);
        \draw[arrow] (brk) -- node[below, font=\scriptsize]{subscribe} (dash);
    \end{tikzpicture}}
    \caption{รูปแบบ publish/subscribe ของ MQTT ในระบบดูแล: gateway ส่งข้อมูลผ่าน broker ไปยัง server และ dashboard}
    \label{fig:ch2_mqtt_pubsub}
\end{figure}


เซ็นเซอร์ IMU และ wearable sensor เป็นแหล่งข้อมูลสำคัญในงานดูแลผู้ใช้รถเข็น IMU มักประกอบด้วย accelerometer และ gyroscope สำหรับวัดความเร่งและอัตราการหมุน ส่วน wearable sensor เช่นอุปกรณ์วัดแสงสะท้อนจากผิวหนังด้วยหลักการ Photoplethysmography (PPG) สามารถใช้ติดตาม heart rate, pulse-to-pulse interval หรือข้อมูลสัญญาณชีพที่สัมพันธ์กับภาวะร่างกายได้ \cite{PolarBLESDK,PolarVeritySenseWhitepaper} อย่างไรก็ตาม ข้อมูลเหล่านี้ควรถูกใช้เป็นข้อมูลประกอบการประเมินเบื้องต้น ไม่ใช่การวินิจฉัยทางการแพทย์

ในระบบที่ต้องประมวลผลภายในพื้นที่ทดลอง แนวคิด edge หรือ local server มีบทบาทในการลดการพึ่งพา cloud ภายนอก ทำให้ข้อมูลส่วนบุคคลไม่จำเป็นต้องออกนอกระบบ และช่วยให้ dashboard หรือระบบแจ้งเตือนทำงานได้แม้ internet ภายนอกไม่เสถียร โครงสร้างทั่วไปมักประกอบด้วย gateway ที่รับข้อมูลจากอุปกรณ์ภาคสนาม server ที่ตรวจสอบและจัดเก็บข้อมูล ฐานข้อมูลสำหรับ telemetry และ frontend/dashboard สำหรับแสดงผล

\begin{figure}[htbp]
    \centering
    \includegraphics[width=0.72\textwidth]{chapter2/homeasisitant.png}
    \caption{ตัวอย่างหน้าจอ dashboard ของระบบ smart environment ที่รวมสถานะจากอุปกรณ์ IoT หลายชนิด}
    \label{fig:ch2_smart_environment_dashboard}
\end{figure}

\begin{table}[htbp]
    \centering
    \caption{บทบาทของเทคโนโลยี IoT และ sensor ในระบบดูแล}
    \label{tab:ch2_iot_sensor_roles}
    \small
    \begin{tabular}{>{\raggedright\arraybackslash}p{0.23\textwidth}
                    >{\raggedright\arraybackslash}p{0.32\textwidth}
                    >{\raggedright\arraybackslash}p{0.34\textwidth}}
        \toprule
        เทคโนโลยี & หลักการสำคัญ & เหตุผลที่เกี่ยวข้องกับระบบดูแล \\ \midrule
        BLE & สื่อสารระยะใกล้ ใช้พลังงานต่ำ รองรับ advertising และ GATT & ใช้เชื่อม sensor กับ mobile gateway และใช้ RSSI เป็นข้อมูลตำแหน่งภายในอาคาร \\ \midrule
        MQTT & publish/subscribe ผ่าน broker & ส่ง telemetry หลายชนิดไปยัง server โดยไม่ผูกอุปกรณ์เข้ากับ service โดยตรง \\ \midrule
        IMU & วัดความเร่งและอัตราการหมุนหลายแกน & ใช้วิเคราะห์การเคลื่อนไหวของรถเข็นหรือกิจกรรมพื้นฐาน \\ \midrule
        Wearable sensor & วัดข้อมูลชีวสัญญาณ เช่น HR, RR interval และ PPG & เพิ่มข้อมูลสถานะร่างกายเพื่อประกอบการเฝ้าระวัง \\ \midrule
        Edge/local server & ประมวลผลและจัดเก็บข้อมูลในพื้นที่ & ลด latency และลดการพึ่งพาระบบ cloud ภายนอก \\ \bottomrule
    \end{tabular}
\end{table}

\section{อุปกรณ์ฮาร์ดแวร์หลักที่เกี่ยวข้อง}
\label{sec:ch2_hardware}

M5StickC Plus2 เป็นบอร์ดไมโครคอนโทรลเลอร์ขนาดเล็กที่มี ESP32, หน้าจอ, ปุ่ม, แบตเตอรี่ และ IMU ในตัว จึงเหมาะกับการทดลองเป็น sensor ติดรถเข็นเพื่ออ่านค่าการเคลื่อนไหวและสถานะแบตเตอรี่ ข้อดีของอุปกรณ์ลักษณะนี้คือมีขนาดเล็ก ติดตั้งง่าย และสามารถส่งข้อมูลผ่าน BLE ไปยังอุปกรณ์กลางได้ ในบริบทของงานนี้ M5StickC Plus2 ถูกพิจารณาในฐานะ wheelchair sensor ที่อ่านข้อมูล IMU/battery และส่ง payload ไปยัง mobile gateway ไม่ใช่อุปกรณ์ที่ต้องทำหน้าที่เป็น gateway หลักของระบบทั้งหมด \cite{M5StickCPlus2Docs}

LilyGO T-SimCam เป็นบอร์ดที่มีความสามารถด้าน cellular และ camera module จึงเคยเหมาะกับแนวคิด node หรือ camera node ในต้นแบบช่วงแรก โดยเฉพาะการทดลองโครงข่ายที่ต้องส่งข้อมูลหรือภาพผ่าน node ภาคสนาม อย่างไรก็ตาม อุปกรณ์ประเภทนี้มีความซับซ้อนด้านการติดตั้ง พลังงาน และการจัดเส้นทางข้อมูลมากกว่า sensor ขนาดเล็ก ในระบบสุดท้ายจึงควรถูกอธิบายในฐานะอุปกรณ์ที่เคยใช้ในต้นแบบหรือแนวคิด node เดิม ไม่ใช่แกนหลักของ architecture ปัจจุบัน \cite{LilygoTSimCamDocs}

Polar Verity Sense เป็นอุปกรณ์สวมใส่แบบ optical armband สำหรับอ่านข้อมูล heart rate และข้อมูลที่เกี่ยวข้องกับ PPG ผ่านการเชื่อมต่อ BLE และ Polar SDK อุปกรณ์นี้มีบทบาทในการเพิ่มข้อมูลด้านสัญญาณชีพให้ระบบดูแล เช่น HR, RR interval หรือ PPG stream โดยการอ่านค่าระดับลึกต้องอาศัย SDK ที่ผู้ผลิตเตรียมไว้ ข้อมูลจากอุปกรณ์สวมใส่จึงช่วยเสริมข้อมูลการเคลื่อนไหวจากรถเข็น แต่ยังต้องถูกตีความอย่างระมัดระวังในฐานะข้อมูลประกอบ ไม่ใช่ผลวินิจฉัย \cite{Polar_VeritySense_Spec,PolarBLESDK}

Raspberry Pi 5 RAM 8 GB เป็น single-board computer ที่มีพลังประมวลผลเพียงพอสำหรับบทบาท local/edge server ในระบบต้นแบบ เช่น API service, database, MQTT broker, dashboard และ AI runtime ขนาดเล็ก การใช้เครื่องลักษณะนี้ช่วยให้ระบบทดลองสามารถทำงานในพื้นที่ได้โดยไม่ต้องพึ่ง cloud ทั้งหมด และยังเปิดโอกาสให้ทดสอบการ deploy ด้วย Docker, การจัดเก็บข้อมูลใน PostgreSQL และการเรียกใช้ AI/MCP ภายในระบบเดียวกัน \cite{RaspberryPi5Specs}

\begin{figure}[htbp]
    \centering
    \begin{subfigure}[t]{0.45\textwidth}
        \centering
        \includegraphics[width=0.72\linewidth]{chapter2/ch2-m5stick-product.png}
        \caption{M5StickC Plus2}
        \label{fig:ch2_m5stick_product}
    \end{subfigure}
    \hfill
    \begin{subfigure}[t]{0.45\textwidth}
        \centering
        \includegraphics[width=0.72\linewidth]{chapter2/ch2-tsimcam-product.png}
        \caption{LilyGO T-SimCam}
        \label{fig:ch2_tsimcam_product}
    \end{subfigure}

    \vspace{0.4cm}
    \begin{subfigure}[t]{0.45\textwidth}
        \centering
        \includegraphics[width=0.72\linewidth]{chapter2/ch2-polar-verity-sense.png}
        \caption{Polar Verity Sense}
        \label{fig:ch2_polar_product}
    \end{subfigure}
    \hfill
    \begin{subfigure}[t]{0.45\textwidth}
        \centering
        \includegraphics[width=0.78\linewidth]{chapter2/pi5.jpg}
        \caption{Raspberry Pi 5}
        \label{fig:ch2_pi5}
    \end{subfigure}
    \caption{อุปกรณ์ฮาร์ดแวร์หลักที่เกี่ยวข้องกับระบบต้นแบบ}
    \label{fig:ch2_hardware_devices}
\end{figure}

\begin{table}[htbp]
    \centering
    \caption{สเปคสำคัญของอุปกรณ์ฮาร์ดแวร์หลัก}
    \label{tab:ch2_hardware_specs}
    \small
    \begin{tabular}{>{\raggedright\arraybackslash}p{0.22\textwidth}
                    >{\raggedright\arraybackslash}p{0.20\textwidth}
                    >{\raggedright\arraybackslash}p{0.24\textwidth}
                    >{\raggedright\arraybackslash}p{0.24\textwidth}}
        \toprule
        คุณสมบัติ & M5StickC Plus2 & Polar Verity Sense & Raspberry Pi 5 \\ \midrule
        MCU / CPU & ESP32-PICO-V3 & Nordic nRF5340 & BCM2712 Cortex-A76 \\ \midrule
        Sensor & IMU 6 แกน (MPU-6886) & PPG (HR, RR interval) & -- \\ \midrule
        การเชื่อมต่อ & BLE 4.2, WiFi & BLE, ANT+ & WiFi 5, GbE, BLE 5.0 \\ \midrule
        แบตเตอรี่ & 200 mAh LiPo & 110 mAh (สวมใส่ 20+ ชม.) & ต้องใช้ไฟภายนอก 5V/5A \\ \midrule
        RAM & 520 KB SRAM & -- & 8 GB LPDDR4X \\ \midrule
        บทบาทในระบบ & Wheelchair sensor & Wearable sensor & Edge / local server \\ \bottomrule
    \end{tabular}
\end{table}

นอกจากอุปกรณ์หลักข้างต้น ระบบดูแลยังอาจมีอุปกรณ์รอง เช่น BLE beacon หรือ node อ้างอิงสัญญาณ โทรศัพท์มือถือที่ทำหน้าที่ gateway เครื่องชาร์จ แหล่งจ่ายไฟสำรอง และอุปกรณ์เครือข่ายภายในอาคาร อุปกรณ์เหล่านี้มีผลต่อความเสถียรของระบบ แต่ในเชิงทฤษฎีสามารถจัดเป็นส่วนสนับสนุนของ sensor layer, network layer และ edge layer ตามบทบาทที่กล่าวไว้ก่อนหน้า

\FloatBarrier

\section{เครื่องมือซอฟต์แวร์และไลบรารีที่ใช้พัฒนา}
\label{sec:ch2_software_tools}

นอกจากฮาร์ดแวร์แล้ว ระบบต้นแบบยังอาศัยเครื่องมือซอฟต์แวร์หลายชนิดที่ทำงานร่วมกัน ตารางที่~\ref{tab:ch2_software_tools} สรุปเครื่องมือหลักที่เกี่ยวข้องกับการพัฒนาระบบ โดยแต่ละเครื่องมือถูกเลือกจากเหตุผลเชิงเทคนิคและข้อจำกัดของ prototype

\begin{table}[htbp]
    \centering
    \caption{เครื่องมือซอฟต์แวร์และไลบรารีหลักที่ใช้ในระบบต้นแบบ}
    \label{tab:ch2_software_tools}
    \small
    \begin{tabular}{>{\raggedright\arraybackslash}p{0.18\textwidth}
                    >{\raggedright\arraybackslash}p{0.14\textwidth}
                    >{\raggedright\arraybackslash}p{0.30\textwidth}
                    >{\raggedright\arraybackslash}p{0.28\textwidth}}
        \toprule
        เครื่องมือ & ประเภท & หน้าที่ & เหตุผลที่เลือกใช้ \\ \midrule
        Flutter & Mobile SDK & พัฒนา mobile gateway app สำหรับ Android & รองรับ BLE plugin, cross-platform, UI สำหรับ caregiver \\ \midrule
        Next.js & Web framework & พัฒนา dashboard สำหรับ admin และ caregiver & Server-side rendering, real-time update \\ \midrule
        FastAPI & Backend API & RESTful API, WebSocket, business logic & Async, type-safe, รองรับ Python AI ecosystem \\ \midrule
        PostgreSQL & Database & จัดเก็บ telemetry, user, alert ด้วย JSONB & รองรับ JSON payload ยืดหยุ่น \\ \midrule
        Docker & Containerization & Deploy service ทั้งหมดบน Raspberry Pi 5 & แยก environment, ทำซ้ำได้ \\ \midrule
        Ollama + Gemma 3 4B & AI runtime & Local LLM สำหรับ AI assistant & ทำงานบน edge ได้, ไม่ต้องพึ่ง cloud \\ \midrule
        FastMCP & MCP server & กำหนด tool contract สำหรับ AI เรียกใช้ & สอดคล้องกับ MCP specification \\ \midrule
        PlatformIO & Embedded IDE & พัฒนา firmware สำหรับ M5StickC Plus2 & รองรับ ESP32, จัดการ library อัตโนมัติ \\ \bottomrule
    \end{tabular}
\end{table}

Flutter เป็น SDK สำหรับพัฒนาแอปพลิเคชันมือถือแบบ cross-platform ที่ใช้ภาษา Dart สามารถเชื่อมต่อ BLE ผ่าน plugin และสร้าง UI สำหรับแสดงสถานะ sensor ได้ ในระบบนี้ Flutter ถูกใช้พัฒนา mobile gateway ที่รับข้อมูลจาก sensor ติดรถเข็นและ wearable sensor ผ่าน BLE แล้วส่งต่อไปยัง server ผ่าน MQTT หรือ REST API \cite{FlutterDocs}

Next.js เป็น React-based web framework ที่รองรับ server-side rendering และ API routes ใช้พัฒนา dashboard สำหรับ admin และ caregiver ให้สามารถดูข้อมูลผู้พักอาศัย สถานะ sensor ตำแหน่งภายในอาคาร และประวัติการแจ้งเตือน \cite{NextJS2024}

FastAPI เป็น Python web framework แบบ asynchronous ที่รองรับ type hint และ auto-documentation ทำหน้าที่เป็น backend หลักของระบบ รับ telemetry จาก gateway, ประมวลผล business logic, จัดเก็บข้อมูลใน PostgreSQL และให้บริการ WebSocket สำหรับ real-time update \cite{FastAPI2024}

PostgreSQL เป็นระบบฐานข้อมูลเชิงสัมพันธ์แบบ open source ที่รองรับ JSONB ทำให้สามารถจัดเก็บ telemetry payload ที่มีโครงสร้างยืดหยุ่นได้โดยไม่ต้องเปลี่ยน schema ทุกครั้งที่เพิ่มข้อมูลใหม่ \cite{PostgreSQL2024}

Docker เป็นเครื่องมือ containerization ที่ช่วยให้ deploy service ทุกส่วนบน Raspberry Pi 5 ด้วย docker-compose ทำให้แต่ละ service แยก environment และทำซ้ำได้ \cite{DockerCompose2024}

Ollama เป็น runtime สำหรับเรียกใช้ Large Language Model บน local machine ในระบบนี้ใช้ร่วมกับ Gemma 3 4B เพื่อให้ AI assistant ทำงานบน edge server ได้โดยไม่ต้องส่งข้อมูลไปยัง cloud ภายนอก \cite{Ollama2024}

FastMCP เป็น Python library สำหรับสร้าง MCP server ตาม Model Context Protocol specification ทำให้ AI client สามารถเรียกใช้เครื่องมือภายในระบบผ่าน tool contract ที่มีการกำหนด input/output schema ชัดเจน \cite{ModelContextProtocol}

\FloatBarrier

\section{การระบุตำแหน่งภายในอาคารด้วย BLE RSSI}
\label{sec:ch2_ble_rssi_localization}

การระบุตำแหน่งภายในอาคารมีความแตกต่างจากการระบุตำแหน่งกลางแจ้ง เนื่องจาก GPS ทำงานได้ไม่ดีในอาคารที่มีผนัง เพดาน และสัญญาณสะท้อนจำนวนมาก ระบบภายในอาคารจึงมักใช้สัญญาณวิทยุจาก WiFi, BLE, UWB หรือ RFID แทน สำหรับระบบดูแลผู้สูงอายุ BLE เป็นตัวเลือกที่เหมาะกับต้นแบบเพราะอุปกรณ์มีราคาต่ำ ใช้พลังงานน้อย และสามารถวัด RSSI ได้จากอุปกรณ์รับสัญญาณทั่วไป

แนวคิดพื้นฐานของ RSSI คือความแรงของสัญญาณที่ตัวรับวัดได้จะลดลงเมื่อระยะห่างจากตัวส่งเพิ่มขึ้น แบบจำลองที่ใช้บ่อยคือ log-distance path loss model ดังสมการ~(\ref{eq:ch2_rssi_distance}) \cite{Kaemarungsi2012RSSIAnalysis}

\begin{equation}
    d = d_0 \times 10^{\frac{P(d_0)-P(d)}{10n}}
    \label{eq:ch2_rssi_distance}
\end{equation}

เมื่อ $d$ คือระยะที่ต้องการประมาณ, $d_0$ คือระยะอ้างอิง, $P(d_0)$ คือ RSSI ที่ระยะอ้างอิง, $P(d)$ คือ RSSI ที่วัดได้ และ $n$ คือ path loss exponent ของสภาพแวดล้อม แม้สมการนี้ช่วยอธิบายความสัมพันธ์ระหว่าง RSSI กับระยะทางได้ แต่ในอาคารจริง RSSI มีความผันผวนสูงจาก multipath, สิ่งกีดขวาง และทิศทางของเสาอากาศ ดังนั้นการคำนวณระยะตรง ๆ แล้วหาจุดตัดแบบ trilateration จึงไม่ใช่แกนหลักที่เหมาะกับระบบนี้

วิธี fingerprinting แก้ปัญหาดังกล่าวโดยแบ่งงานเป็นสองช่วง ช่วงแรกคือ offline phase ที่วัดค่า RSSI จากตำแหน่งอ้างอิงหลายจุดและบันทึกเป็น radio map ช่วงที่สองคือ online phase ที่นำ RSSI ที่วัดได้จริงมาเปรียบเทียบกับ radio map เพื่อหาตำแหน่งที่มีรูปแบบสัญญาณใกล้เคียงที่สุด แนวคิดนี้ถูกใช้ในงาน indoor localization อย่างแพร่หลายตั้งแต่งานกลุ่ม RADAR และต่อยอดมาในระบบ BLE RSSI หลายรูปแบบ \cite{Bahl2000RADAR}

K-Nearest Neighbors (KNN) เป็นวิธีพื้นฐานที่ใช้กับ fingerprinting ได้ตรงไปตรงมา โดยแทนค่า RSSI ที่วัดได้เป็นเวกเตอร์ $\mathbf{x}$ และเปรียบเทียบกับ fingerprint ที่เก็บไว้เป็น $\mathbf{r}_j$ ระยะห่างระหว่างเวกเตอร์คำนวณได้ตามสมการ~(\ref{eq:ch2_knn_rssi_distance}) \cite{Cover1967NearestNeighbor}

\begin{equation}
    D_j = \sqrt{\sum_{i=1}^{m}(x_i-r_{j,i})^2}
    \label{eq:ch2_knn_rssi_distance}
\end{equation}

เมื่อ $m$ คือจำนวน beacon หรือ node อ้างอิง, $x_i$ คือ RSSI ที่วัดได้จาก node ที่ $i$ และ $r_{j,i}$ คือค่า RSSI ของ fingerprint ลำดับที่ $j$ จากนั้นเลือก fingerprint ที่มีระยะน้อยที่สุดจำนวน $K$ จุด เพื่อนำมาประเมินห้องหรือพิกัด การใช้ distance-weighted KNN ช่วยให้น้ำหนักของจุดที่ใกล้กว่าใน feature space มากกว่าจุดที่อยู่ไกลกว่า ดังสมการ~(\ref{eq:ch2_weighted_knn})

\begin{equation}
    w_j = \frac{1}{D_j+\epsilon}, \qquad
    \hat{c} = \operatorname*{arg\,max}_{c} \sum_{j \in N_K} w_j \cdot I(c_j=c)
    \label{eq:ch2_weighted_knn}
\end{equation}

เมื่อ $N_K$ คือชุด fingerprint ใกล้เคียงที่สุดจำนวน $K$ จุด, $c_j$ คือคลาสหรือห้องของ fingerprint จุดที่ $j$, $I(\cdot)$ คือ indicator function และ $\epsilon$ เป็นค่าขนาดเล็กเพื่อป้องกันการหารด้วยศูนย์

ในกรณีที่ข้อมูล fingerprint ไม่ครบหรือค่า RSSI ไม่เหมาะกับการคำนวณ KNN ระบบสามารถใช้ max-RSSI เป็น baseline หรือ fallback ได้ แนวคิดคือเลือก node ที่มีค่า RSSI สูงที่สุดแล้วใช้ตำแหน่งหรือห้องของ node นั้นเป็นค่าประมาณ ดังสมการ~(\ref{eq:ch2_max_rssi})

\begin{equation}
    i^{*} = \operatorname*{arg\,max}_{i} P_i, \qquad
    \hat{c}_{\mathrm{max}} = c_{i^{*}}
    \label{eq:ch2_max_rssi}
\end{equation}

วิธี max-RSSI มีความเรียบง่ายและเหมาะเป็น fallback แต่โดยทั่วไปจะไวต่อสัญญาณสะท้อนหรือการบังสัญญาณมากกว่า fingerprinting ดังนั้นในระบบต้นแบบจึงควรใช้ KNN เป็นวิธีหลักเมื่อมีข้อมูลเพียงพอ และใช้ max-RSSI ในกรณีที่ข้อมูลไม่ครบหรือโมเดลไม่สามารถประเมินได้อย่างมั่นใจ

\begin{figure}[htbp]
    \centering
    \resizebox{0.94\textwidth}{!}{%
    \begin{tikzpicture}[
        font=\small,
        box/.style={draw, thick, align=center, text width=3.0cm, minimum height=1.0cm},
        arrow/.style={-{Latex[length=2mm]}, thick}
    ]
        \node[box, fill=gray!10] (collect) {วัด RSSI\\ตามจุดอ้างอิง};
        \node[box, fill=blue!8, right=0.75cm of collect] (map) {สร้าง radio map\\fingerprint};
        \node[box, fill=green!8, right=0.75cm of map] (online) {รับ RSSI ใหม่\\จาก gateway};
        \node[box, fill=orange!10, right=0.75cm of online] (knn) {KNN / weighted vote\\หรือ max-RSSI};
        \node[box, fill=purple!8, right=0.75cm of knn] (room) {ผลลัพธ์\\ห้องหรือพื้นที่};
        \draw[arrow] (collect) -- (map);
        \draw[arrow] (map) -- (online);
        \draw[arrow] (online) -- (knn);
        \draw[arrow] (knn) -- (room);
    \end{tikzpicture}}
    \caption{กระบวนการระบุตำแหน่งภายในอาคารด้วย BLE RSSI และ fingerprinting}
    \label{fig:ch2_ble_fingerprinting_pipeline}
\end{figure}

\section{การประมวลผลข้อมูล IMU และการเคลื่อนไหว}
\label{sec:ch2_imu_processing}

IMU เป็นเซ็นเซอร์ที่ใช้วัดการเคลื่อนไหวเชิงกล โดย accelerometer วัดความเร่งเชิงเส้นในแกน $x$, $y$, $z$ และ gyroscope วัดอัตราการหมุนรอบแกนต่าง ๆ ข้อมูลทั้งสองชนิดสามารถใช้ตรวจจับการเปลี่ยนแปลงสถานะของรถเข็นหรือกิจกรรมของผู้ใช้ได้ เช่น การเริ่มเคลื่อนที่ การหยุดนิ่ง การเคลื่อนที่ต่อเนื่อง หรือการสั่นผิดปกติ หลักการประมวลผลข้อมูล IMU จึงมักเริ่มจากการลด noise การคำนวณ feature ในช่วงเวลา และการสรุปสถานะจาก feature เหล่านั้น \cite{Lara2013WearableHAR}

สำหรับ accelerometer สามารถคำนวณ magnitude ของความเร่งรวมได้ตามสมการ~(\ref{eq:ch2_acc_magnitude}) เพื่อรวมข้อมูลสามแกนให้เป็นค่าเดียวที่อ่านง่ายขึ้น

\begin{equation}
    |\mathbf{a}_t| = \sqrt{a_{x,t}^{2}+a_{y,t}^{2}+a_{z,t}^{2}}
    \label{eq:ch2_acc_magnitude}
\end{equation}

ข้อมูลที่วัดต่อเนื่องมักถูกแบ่งเป็น sliding window ความยาว $W$ ตัวอย่าง แล้วคำนวณค่าสถิติ เช่น ค่าเฉลี่ยและส่วนเบี่ยงเบนมาตรฐาน ดังสมการ~(\ref{eq:ch2_window_stats})

\begin{equation}
    \mu = \frac{1}{W}\sum_{t=1}^{W}x_t, \qquad
    \sigma = \sqrt{\frac{1}{W-1}\sum_{t=1}^{W}(x_t-\mu)^2}
    \label{eq:ch2_window_stats}
\end{equation}

อีก feature หนึ่งที่ใช้บ่อยคือจำนวนครั้งที่สัญญาณเปลี่ยนเครื่องหมายหรือ zero-crossing ซึ่งสะท้อนความถี่ของการสั่นหรือการเปลี่ยนทิศทางในช่วงเวลานั้น ดังสมการ~(\ref{eq:ch2_zero_crossing})

\begin{equation}
    ZC = \sum_{t=2}^{W} I(x_t x_{t-1} < 0)
    \label{eq:ch2_zero_crossing}
\end{equation}

ในงานที่นำ IMU ไปติดตั้งกับล้อรถเข็น ข้อมูลจาก gyroscope สามารถใช้ประมาณความเร็วและระยะทางของล้อได้ หากให้ $\omega_t$ เป็นอัตราการหมุนที่ผ่านการกรองแล้วในหน่วย rad/s และ $R$ เป็นรัศมีล้อ ความเร็วเชิงเส้นโดยประมาณคำนวณได้จากสมการ~(\ref{eq:ch2_wheel_velocity})

\begin{equation}
    v_t = R\omega_t
    \label{eq:ch2_wheel_velocity}
\end{equation}

จากนั้นสามารถประมาณความเร่งและระยะทางสะสมได้ตามสมการ~(\ref{eq:ch2_wheel_acc_distance})

\begin{equation}
    a_t = \frac{v_t-v_{t-1}}{\Delta t}, \qquad
    s_t = s_{t-1} + v_t \Delta t
    \label{eq:ch2_wheel_acc_distance}
\end{equation}

ในทางปฏิบัติ ข้อมูล gyroscope ต้องผ่านการชดเชย offset และตัวกรอง เช่น Exponential Moving Average (EMA) เพื่อให้ค่าที่ใช้คำนวณความเร็วไม่แกว่งมากเกินไป สมการทั่วไปของ EMA แสดงในสมการ~(\ref{eq:ch2_ema})

\begin{equation}
    \hat{\omega}_{t} = \alpha \omega_t + (1-\alpha)\hat{\omega}_{t-1}
    \label{eq:ch2_ema}
\end{equation}

เมื่อ $\alpha$ เป็นค่าปรับแต่งของตัวกรอง และ $\hat{\omega}_{t}$ คือสัญญาณหลังกรอง นอกจากนี้มักกำหนด deadband เพื่อตัด drift ขนาดเล็กขณะรถเข็นหยุดนิ่ง ดังสมการ~(\ref{eq:ch2_deadband})

\begin{equation}
    \omega^{'}_t =
    \begin{cases}
        0, & |\hat{\omega}_t| < \tau \\
        \hat{\omega}_t, & |\hat{\omega}_t| \geq \tau
    \end{cases}
    \label{eq:ch2_deadband}
\end{equation}

โดย $\tau$ คือ threshold ของ deadband ที่ต้องปรับให้เหมาะกับอุปกรณ์และตำแหน่งติดตั้ง การนำสูตรเหล่านี้ไปใช้จึงต้องประเมินร่วมกับการทดสอบจริง เพราะ IMU มี drift, noise, การสั่นจากพื้น และความคลาดเคลื่อนจากการติดตั้ง

\begin{table}[htbp]
    \centering
    \caption{ตัวอย่าง feature จาก IMU ที่ใช้กับการวิเคราะห์การเคลื่อนไหว}
    \label{tab:ch2_imu_features}
    \small
    \begin{tabular}{>{\raggedright\arraybackslash}p{0.25\textwidth}
                    >{\raggedright\arraybackslash}p{0.30\textwidth}
                    >{\raggedright\arraybackslash}p{0.34\textwidth}}
        \toprule
        Feature & ข้อมูลที่ใช้ & ความหมายเชิงพฤติกรรม \\ \midrule
        Acceleration magnitude & $a_x,a_y,a_z$ & ระดับแรงสั่นหรือการเปลี่ยนการเคลื่อนไหวโดยรวม \\ \midrule
        Mean และ standard deviation & ค่าใน sliding window & ความเข้มข้นและความแปรปรวนของการเคลื่อนไหว \\ \midrule
        Zero-crossing & สัญญาณหลังกรองในแต่ละแกน & ความถี่ของการสลับทิศทางหรือการสั่น \\ \midrule
        Velocity และ distance & Gyroscope และรัศมีล้อ & ความเร็วและระยะทางโดยประมาณของรถเข็น \\ \bottomrule
    \end{tabular}
\end{table}

\section{Machine Learning สำหรับจำแนกพฤติกรรมหรือความเสี่ยง}
\label{sec:ch2_ml_behavior}

Machine Learning ในระบบดูแลผู้สูงอายุทำหน้าที่แปลงข้อมูล sensor ที่มีความถี่สูงและอ่านยากให้เป็นสถานะหรือความเสี่ยงที่ผู้ดูแลเข้าใจได้ง่ายขึ้น ตัวอย่างเช่น แปลงข้อมูล IMU เป็นสถานะกำลังเคลื่อนที่ หยุดนิ่ง หรือมีการสั่นผิดปกติ หรือรวมข้อมูลการเคลื่อนไหวกับข้อมูลสัญญาณชีพเพื่อช่วยจัดลำดับความสำคัญของเหตุการณ์ ขั้นตอนสำคัญคือ feature engineering ซึ่งนำข้อมูลดิบมาสรุปเป็นตัวแปร เช่น mean, standard deviation, magnitude, zero-crossing, velocity, distance, heart rate และ RR interval

XGBoost เป็นอัลกอริทึม gradient tree boosting ที่ได้รับความนิยมเพราะรองรับ feature หลายชนิด จัดการความสัมพันธ์ไม่เชิงเส้นได้ดี และให้ผลลัพธ์ที่ตีความเชิง feature importance ได้ระดับหนึ่ง โมเดลประกอบด้วย decision tree หลายต้น โดยผลทำนายเกิดจากผลรวมของต้นไม้ทั้งหมด ดังสมการ~(\ref{eq:ch2_xgboost_model}) \cite{Chen2016XGBoost}

\begin{equation}
    \hat{y}_i = \sum_{m=1}^{M} f_m(\mathbf{x}_i), \qquad f_m \in \mathcal{F}
    \label{eq:ch2_xgboost_model}
\end{equation}

เมื่อ $\mathbf{x}_i$ คือ feature vector ของตัวอย่างที่ $i$, $f_m$ คือ decision tree ต้นที่ $m$ และ $M$ คือจำนวนต้นไม้ทั้งหมด โดยการฝึกโมเดลจะพยายามลด loss ของข้อมูลฝึกพร้อมกับ regularization เพื่อลด overfitting ดังสมการ~(\ref{eq:ch2_xgboost_obj})

\begin{equation}
    \mathcal{L} = \sum_{i=1}^{n} l(y_i,\hat{y}_i) + \sum_{m=1}^{M}\Omega(f_m)
    \label{eq:ch2_xgboost_obj}
\end{equation}

หากโมเดลให้ผลลัพธ์เป็น score ของแต่ละคลาส สามารถแปลงเป็นความน่าจะเป็นด้วย softmax ได้ตามสมการ~(\ref{eq:ch2_softmax})

\begin{equation}
    p(y=c|\mathbf{x}) = \frac{\exp(z_c)}{\sum_{q=1}^{C}\exp(z_q)}
    \label{eq:ch2_softmax}
\end{equation}

การประเมินโมเดล classification ควรพิจารณามากกว่า accuracy เพียงค่าเดียว เพราะข้อมูลพฤติกรรมหรือความเสี่ยงอาจไม่สมดุลระหว่างคลาส ค่า precision, recall และ F1-score จึงใช้ประกอบกันตามสมการ~(\ref{eq:ch2_classification_metrics}) \cite{Sokolova2009ClassificationMetrics}

\begin{equation}
    Accuracy = \frac{TP+TN}{TP+TN+FP+FN}, \quad
    Precision = \frac{TP}{TP+FP}, \quad
    Recall = \frac{TP}{TP+FN}
    \label{eq:ch2_classification_metrics}
\end{equation}

\begin{equation}
    F1 = 2 \times \frac{Precision \times Recall}{Precision + Recall}
    \label{eq:ch2_f1}
\end{equation}

สำหรับระบบดูแล ความผิดพลาดของโมเดลมีความหมายเชิงความปลอดภัยต่างกัน False negative อาจทำให้ระบบพลาดเหตุการณ์สำคัญ ขณะที่ false positive อาจเพิ่มภาระผู้ดูแลจาก alert ที่ไม่จำเป็น ดังนั้นการใช้โมเดลควรถูกออกแบบให้เป็นตัวช่วยคัดกรองเบื้องต้นและให้มนุษย์ตรวจสอบ ไม่ใช่กลไกตัดสินใจแทนผู้ดูแลทั้งหมด

\section{Retrieval-Augmented Generation และ Retrieval Logic}
\label{sec:ch2_rag_retrieval_logic}

Retrieval-Augmented Generation (RAG) เป็นแนวคิดที่ให้โมเดลภาษาใช้ข้อมูลจากแหล่งความรู้ภายนอกก่อนสร้างคำตอบ แทนที่จะตอบจากความรู้ที่อยู่ในพารามิเตอร์ของโมเดลเพียงอย่างเดียว ขั้นตอนทั่วไปคือรับ query, ค้นเอกสารหรือข้อมูลที่เกี่ยวข้อง, จัดแพ็กบริบทให้กระชับ และส่งบริบทนั้นให้ LLM สังเคราะห์คำตอบ วิธีนี้ช่วยลดปัญหาการตอบจากข้อมูลที่ไม่อยู่ในบริบท และทำให้ระบบสามารถอ้างอิงข้อมูลที่เปลี่ยนแปลงได้ง่ายขึ้น \cite{Lewis2020RAG}

ในเชิงทฤษฎี RAG-Sequence สามารถเขียนแบบย่อได้ตามสมการ~(\ref{eq:ch2_rag_sequence}) โดย $x$ คือ query, $z$ คือเอกสารหรือบริบทที่ retrieve ได้ และ $y$ คือคำตอบที่โมเดลสร้าง

\begin{equation}
    p_{\mathrm{RAG\text{-}Seq}}(y|x)
    \approx
    \sum_{z \in top\text{-}k(p_{\eta}(\cdot|x))}
    p_{\eta}(z|x)
    \prod_{i=1}^{N} p_{\theta}(y_i|x,z,y_{<i})
    \label{eq:ch2_rag_sequence}
\end{equation}

เมื่อ $p_{\eta}(z|x)$ คือความน่าจะเป็นของ retriever ที่เลือกบริบท $z$ จาก query และ $p_{\theta}$ คือ generator หรือ LLM ที่สร้างคำตอบจาก query และบริบทที่ได้ จุดสำคัญคือคุณภาพของคำตอบไม่ได้ขึ้นกับ LLM เพียงอย่างเดียว แต่ขึ้นกับการเลือกบริบท การจัดลำดับข้อมูล และการตัดข้อมูลที่ไม่จำเป็นออกจาก context window ด้วย

\begin{figure}[htbp]
    \centering
    \resizebox{0.95\textwidth}{!}{%
    \begin{tikzpicture}[
        font=\small,
        box/.style={draw, thick, align=center, text width=3.0cm, minimum height=1.05cm},
        arrow/.style={-{Latex[length=2mm]}, thick}
    ]
        \node[box, fill=gray!10] (query) {Query\\หรือ intent};
        \node[box, fill=blue!8, right=0.65cm of query] (plan) {Plan\\เลือกแหล่งข้อมูล};
        \node[box, fill=green!8, right=0.65cm of plan] (retrieve) {Fetch + Rank\\ดึงและจัดอันดับ};
        \node[box, fill=orange!10, right=0.65cm of retrieve] (pack) {Pack\\ลดซ้ำและจัดบริบท};
        \node[box, fill=purple!8, right=0.65cm of pack] (llm) {LLM\\ตอบจากบริบท};
        \draw[arrow] (query) -- (plan);
        \draw[arrow] (plan) -- (retrieve);
        \draw[arrow] (retrieve) -- (pack);
        \draw[arrow] (pack) -- (llm);
    \end{tikzpicture}}
    \caption{แนวคิด Retrieval Logic ในระบบ RAG}
    \label{fig:ch2_rag_pipeline}
\end{figure}

ในระบบที่เกี่ยวข้องกับข้อมูลผู้ป่วยหรือผู้พักอาศัย retrieval logic ควรมีข้อจำกัดชัดเจน เช่น ดึงเฉพาะข้อมูลที่ผู้ใช้มีสิทธิ์เห็น ดึงเฉพาะช่วงเวลาที่เกี่ยวข้อง และไม่ส่งข้อมูลดิบมากเกินจำเป็นเข้าไปใน context window การแยกขั้นตอน retrieval ออกจากขั้นตอน generation ยังช่วยให้ตรวจสอบได้ว่าคำตอบหนึ่งเกิดจากบริบทใด ซึ่งสำคัญต่อความโปร่งใสและการแก้ไขข้อผิดพลาด

\section{Neuro-symbolic AI และ AI Pipeline แบบ Rule + Model}
\label{sec:ch2_neuro_symbolic_pipeline}

Neuro-symbolic AI คือแนวคิดที่รวมความสามารถของ neural model เช่น deep learning หรือ LLM เข้ากับ symbolic component เช่น กฎ ตรรกะ schema knowledge graph หรือ policy ที่ตรวจสอบได้ งานวิจัยในสาขานี้มองว่าการใช้โมเดลเรียนรู้เพียงอย่างเดียวอาจมีข้อจำกัดด้านความโปร่งใสและการควบคุม ขณะที่ระบบ rule-based เพียงอย่างเดียวอาจไม่ยืดหยุ่นพอสำหรับภาษาธรรมชาติและบริบทที่ซับซ้อน การผสมทั้งสองส่วนจึงมีเป้าหมายเพื่อให้ระบบเข้าใจข้อมูลไม่เป็นโครงสร้างได้ดีขึ้น แต่ยังอยู่ภายใต้ข้อจำกัดที่ตรวจสอบได้ \cite{Susskind2021NeuroSymbolic,ColeloughRegli2025NeuroSymbolic}

สำหรับระบบดูแลที่ใช้ LLM และเครื่องมือภายนอก แนวคิดนี้สามารถนำมาใช้ในลักษณะ neuro-symbolic-inspired ได้ กล่าวคือ neural component ทำหน้าที่เข้าใจภาษา สรุปบริบท และสร้างคำตอบภาษาธรรมชาติ ส่วน symbolic component ทำหน้าที่กำหนด intent schema, policy, role permission, tool contract และ execution gate เพื่อป้องกันคำสั่งที่ไม่ผ่านเงื่อนไข ระบบลักษณะนี้ไม่ควร claim ว่าเป็น neuro-symbolic system เต็มรูปแบบ แต่สามารถอธิบายได้ว่าได้รับแรงบันดาลใจจากการแยกหน้าที่ระหว่าง model reasoning และ rule-based control

\begin{table}[htbp]
    \centering
    \caption{การแบ่งบทบาทระหว่าง neural และ symbolic component}
    \label{tab:ch2_neuro_symbolic_roles}
    \small
    \begin{tabular}{>{\raggedright\arraybackslash}p{0.24\textwidth}
                    >{\raggedright\arraybackslash}p{0.33\textwidth}
                    >{\raggedright\arraybackslash}p{0.34\textwidth}}
        \toprule
        ส่วนประกอบ & ตัวอย่างหน้าที่ & ประโยชน์ต่อระบบดูแล \\ \midrule
        Neural & เข้าใจภาษาธรรมชาติ จัดกลุ่มความหมาย สังเคราะห์คำตอบ & รองรับคำถามหรือคำสั่งของผู้ดูแลที่มีรูปแบบหลากหลาย \\ \midrule
        Symbolic & intent rule, schema, policy, permission, tool contract & จำกัดขอบเขตการเข้าถึงข้อมูลและการสั่งงานให้ตรวจสอบได้ \\ \midrule
        Execution gate & ตรวจ payload ก่อนเรียกใช้ tool หรือ MCP server & ลดความเสี่ยงจาก hallucination หรือคำสั่งที่ไม่ตรงสิทธิ์ \\ \bottomrule
    \end{tabular}
\end{table}

\begin{figure}[htbp]
    \centering
    \resizebox{0.95\textwidth}{!}{%
    \begin{tikzpicture}[
        font=\small,
        box/.style={draw, thick, align=center, text width=2.75cm, minimum height=1.05cm},
        arrow/.style={-{Latex[length=2mm]}, thick}
    ]
        \node[box, fill=gray!10] (input) {Query\\จากผู้ดูแล};
        \node[box, fill=blue!8, right=0.55cm of input] (router) {Intent Router\\rule + semantic};
        \node[box, fill=green!8, right=0.55cm of router] (context) {Context\\schema-bound};
        \node[box, fill=orange!10, right=0.55cm of context] (llm) {LLM\\constrained answer};
        \node[box, fill=red!8, right=0.55cm of llm] (gate) {Execution Gate\\policy check};
        \node[box, fill=purple!8, right=0.55cm of gate] (tool) {MCP Tool\\หรือผลลัพธ์};
        \draw[arrow] (input) -- (router);
        \draw[arrow] (router) -- (context);
        \draw[arrow] (context) -- (llm);
        \draw[arrow] (llm) -- (gate);
        \draw[arrow] (gate) -- (tool);
    \end{tikzpicture}}
    \caption{แนวคิด AI pipeline แบบผสม rule และ model สำหรับระบบที่เรียกใช้เครื่องมือ}
    \label{fig:ch2_neuro_symbolic_pipeline}
\end{figure}

Model Context Protocol (MCP) เป็นแนวทางหนึ่งในการกำหนดรูปแบบการเชื่อมต่อระหว่าง AI client กับเครื่องมือหรือ resource ภายนอก โดยช่วยให้ tool มี contract ชัดเจนและสามารถควบคุมการเข้าถึงผ่าน server ได้ \cite{ModelContextProtocol,Anthropic2024MCP} เมื่อใช้ร่วมกับ execution gate ระบบสามารถแยกขั้นตอน ``เสนอคำสั่ง'' ออกจาก ``อนุญาตให้สั่งงานจริง'' ได้ชัดเจน ซึ่งเหมาะกับงานดูแลผู้สูงอายุที่การทำงานผิดพลาดอาจส่งผลต่อความปลอดภัยหรือความเป็นส่วนตัวของข้อมูล

\section{งานวิจัยที่เกี่ยวข้อง}
\label{sec:ch2_related_work}

งานวิจัยด้าน indoor localization ด้วย BLE RSSI และ fingerprinting แสดงให้เห็นว่า RSSI สามารถใช้เป็นสัญญาณประกอบการระบุตำแหน่งภายในอาคารได้ แต่มีความผันผวนสูงและต้องอาศัยวิธีลดผลกระทบจากสภาพแวดล้อม งานของ Bahl และ Padmanabhan ในระบบ RADAR เป็นตัวอย่างสำคัญของการใช้ radio map และ fingerprinting เพื่อหาตำแหน่งภายในอาคาร \cite{Bahl2000RADAR} ส่วนงานวิเคราะห์ RSSI ในอาคารชี้ให้เห็นว่าค่า RSSI ได้รับผลกระทบจากตำแหน่งอุปกรณ์ สิ่งกีดขวาง และการสะท้อนของสัญญาณ \cite{Kaemarungsi2012RSSIAnalysis} ทำให้การใช้ KNN และ fallback เช่น max-RSSI มีความเหมาะสมกับ prototype ที่ต้องรองรับความไม่แน่นอนของข้อมูล

งานวิจัยด้าน Human Activity Recognition (HAR) จาก wearable sensor และ IMU มักใช้กระบวนการแบ่งสัญญาณเป็น window แล้วคำนวณ feature เชิงสถิติ เช่น mean, variance, energy, magnitude และ zero-crossing ก่อนป้อนเข้าสู่โมเดลจำแนกพฤติกรรม \cite{Lara2013WearableHAR} แนวทางนี้สอดคล้องกับการใช้ข้อมูลจาก accelerometer และ gyroscope บนรถเข็น เพราะข้อมูลดิบอ่านตีความยาก แต่เมื่อแปลงเป็น feature แล้วสามารถใช้จำแนกสถานะการเคลื่อนไหวหรือความผิดปกติได้ชัดเจนขึ้น

งานด้าน wearable health และ PPG สนับสนุนการใช้ sensor สวมใส่เพื่อติดตามข้อมูลชีวสัญญาณ เช่น heart rate, RR interval, PPG และข้อมูลการเคลื่อนไหว แต่ข้อจำกัดสำคัญคือคุณภาพสัญญาณขึ้นกับตำแหน่งสวมใส่ การเคลื่อนไหวของร่างกาย และ algorithm ของผู้ผลิต \cite{PolarBLESDK,PolarVeritySenseWhitepaper} ดังนั้นการใช้ข้อมูล wearable ในระบบดูแลจึงควรเป็นข้อมูลประกอบการเฝ้าระวังและแจ้งเตือน ไม่ใช่การวินิจฉัยทางการแพทย์

งานด้าน Retrieval-Augmented Generation แสดงให้เห็นว่าการดึงบริบทจากแหล่งข้อมูลภายนอกก่อนให้โมเดลภาษาตอบ สามารถเพิ่มความเฉพาะเจาะจงและความสอดคล้องของคำตอบในงานที่ต้องใช้ความรู้จำนวนมาก \cite{Lewis2020RAG} สำหรับระบบดูแล แนวคิดนี้นำมาประยุกต์ได้กับการดึงข้อมูลผู้พักอาศัย ข้อมูล sensor ล่าสุด ประวัติการแจ้งเตือน หรือ policy ก่อนสร้างคำตอบให้ผู้ดูแล

งานด้าน neuro-symbolic และ tool-use AI เน้นการผสมความสามารถของโมเดลกับกฎหรือ schema ที่ตรวจสอบได้ เพื่อเพิ่มความน่าเชื่อถือและลดความเสี่ยงจากการตัดสินใจแบบ black box \cite{Susskind2021NeuroSymbolic,ColeloughRegli2025NeuroSymbolic} เมื่อระบบต้องเรียกใช้เครื่องมือผ่าน MCP หรือ API ภายนอก การมี execution gate และ tool contract จึงเป็นกลไกสำคัญที่ทำให้การใช้ AI อยู่ในขอบเขตที่ปลอดภัยกว่า LLM ที่สามารถสั่งงานได้โดยไม่มีชั้นตรวจสอบ

\begin{table}[htbp]
    \centering
    \caption{การจัดกลุ่มงานวิจัยที่เกี่ยวข้องกับระบบต้นแบบ}
    \label{tab:ch2_related_work_groups}
    \small
    \begin{tabular}{>{\raggedright\arraybackslash}p{0.27\textwidth}
                    >{\raggedright\arraybackslash}p{0.34\textwidth}
                    >{\raggedright\arraybackslash}p{0.30\textwidth}}
        \toprule
        กลุ่มงานวิจัย & ประเด็นสำคัญ & การเชื่อมโยงกับงานนี้ \\ \midrule
        Indoor localization / BLE RSSI / KNN & RSSI ผันผวน ต้องใช้ fingerprinting หรือวิธี fallback & ใช้ KNN และ max-RSSI เพื่อประเมินพื้นที่ในอาคาร \\ \midrule
        Human activity recognition / IMU & แปลงสัญญาณเป็น feature ก่อนจำแนกพฤติกรรม & ใช้ feature จาก IMU และข้อมูลการเคลื่อนไหวของรถเข็น \\ \midrule
        Wearable health / PPG & วัด HR, RR interval และ PPG แต่ต้องระวัง noise & ใช้ข้อมูล wearable เป็นบริบทประกอบ ไม่ใช่ผลวินิจฉัย \\ \midrule
        RAG & ดึงบริบทก่อนให้ LLM ตอบ & ลดการตอบนอกข้อมูลและควบคุม context ของผู้พักอาศัย \\ \midrule
        Neuro-symbolic / tool-use AI & ใช้ rule, schema และ policy ควบคู่กับโมเดล & ใช้ execution gate และ MCP tool contract ก่อนสั่งงานจริง \\ \bottomrule
    \end{tabular}
\end{table}

\section{ช่องว่างของงานวิจัยเดิม}
\label{sec:ch2_research_gap}

จากการทบทวนทฤษฎีและงานวิจัยที่เกี่ยวข้อง พบว่างานจำนวนมากมักพัฒนาเฉพาะส่วนใดส่วนหนึ่งของระบบ เช่น งาน indoor localization ที่เน้นตำแหน่งจาก RSSI, งาน HAR ที่เน้นจำแนกกิจกรรมจาก IMU, งาน wearable health ที่เน้นสัญญาณชีพ, งาน dashboard ที่เน้นการแสดงผล หรือระบบ AI ที่เน้นการตอบคำถามและการเรียกใช้เครื่องมือ แต่ในบริบทการดูแลผู้สูงอายุหรือผู้ใช้รถเข็น ผู้ดูแลต้องการภาพรวมที่เชื่อมข้อมูลเหล่านี้เข้าด้วยกันมากกว่าผลลัพธ์แยกส่วน

ช่องว่างสำคัญคือการรวม wheelchair sensor, wearable sensor, mobile gateway, indoor localization, dashboard ตามบทบาทผู้ใช้ และ AI/MCP tool pipeline ไว้ใน prototype เดียว โดยยังคงขอบเขตว่าเป็นระบบต้นแบบเพื่อสนับสนุนการดูแล ไม่ใช่เครื่องมือแพทย์ที่ผ่านการรับรอง การออกแบบลักษณะนี้ต้องจัดสมดุลระหว่างความสามารถด้าน sensor, ความเสถียรของการส่งข้อมูล, ความชัดเจนของ dashboard, และความปลอดภัยของ AI ที่อาจเข้าถึงข้อมูลหรือเครื่องมือในระบบ

ดังนั้น บทถัดไปจะนำแนวคิดจากบทนี้ไปออกแบบระบบในรูปแบบปัญหาไปสู่แนวทางแก้ไข โดยอธิบายว่าระบบถูกปรับจากข้อจำกัดของ prototype เดิมอย่างไร เหตุใด mobile gateway จึงกลายเป็นแกนหลักของการส่งข้อมูล และ AI pipeline ถูกแยกเป็นชั้นต่าง ๆ เพื่อควบคุม latency, context, hallucination และการสั่งงานผ่าน MCP ได้อย่างไร
